\scanpage{53}
\uline{NB} Pour bien faire, il faudrait vérifier que la structure de torseur sous $L^{\otimes2}$ qu'on vient d'obtenir sur le faisceau des $q$ (ou $\Delta$) compatibles avec $\underline b$, est aussi celle du § 6.

\uline{10} \uline{Cas mixte}. \deleted{A) \unsure{Soit $S$} Cas où}
\diagram{$\begin{tikzcd}[ampersand replacement=\&,row sep=small,column sep=small]
\deleted{S=}\&\deleted{S_2}\arrow[d]\&\deleted{\cup}\&\deleted{S_D}\arrow[d]\\
\&\deleted{\text{pts où }2\text{ inv.}}\&\&\deleted{\text{pts où }D\text{ inv.}}
\end{tikzcd}$}
\deleted{i.e. où $\Delta$ est étale en tous les pts de $S$ de car. 2.}

\deleted{La donnée de $\Delta$ équivaut alors à celle d'un rev. quadratique \deleted{\unsure{étale}} \inserted{\unsure{induit|radical}} de $S_2$\deleted{\unsure{$D$|$\partial$}} et d'un}

Nous nous donnons comme premier invariant dominant la situation le couple $(L,D)$ d'un Module inv. $L$ et de
\[D\in\Gamma(S,L^{\otimes2})\tag*{(103)}\]
\deleted{et l'hypothèse du titre est prise en ce sens :}

Il faut définir un $\Delta$ compatible avec ces données séparément sur $S_D$ et $S_2$, et \deleted{\unsure{donner|donc}} une donnée de recollement sur $S_2\cap S_D=S_{2D}$.

a) Sur $S_2$, on sait (§ 8) que $\Delta$ est déterminé à isom can. près par $(L,D)$, comme le schéma des solutions de l'équation
\[x^2=D\qquad(x\in\Gamma(-,L))\tag*{(104)}\]
% archive sheet number: 52
